\chapter{12.7 线段的垂直平分线-1}


\section*{三、解答题（共 4 小题）}

\solutionlist[15]

\item 作图题：（不写作法，但必须保留作图痕迹）

如图：某地有两所大学和两条相交的公路，（点 $M, N$ 表示大学，$AO, BO$ 表示公路）。现计划修建一座物资仓库，希望仓库到两所大学的距离相等，到两条公路的距离也相等。你能确定仓库 $P$ 应该建在什么位置吗？在所给的图形中画出你的设计方案。

\rightquestionimage[width=0.4\textwidth]{0.4\textwidth}{12.7-15.png}

\vspace{3cm}

\item 在 $\triangle ABC$ 中，$AD$ 是高，在线段 $DC$ 上取一点 $E$，使 $BD = DE$，已知 $AB + BD = DC$，求证：$E$ 点在线段 $AC$ 的垂直平分线上。

\rightquestionimage[width=0.4\textwidth]{0.4\textwidth}{12.7-16.png}

\vspace{3cm}

\newpage

\item $\triangle ABC$ 中，$DE, FG$ 分别垂直平分边 $AB, AC$，垂足分别为点 $D, G$。

\rightquestionimage[width=0.4\textwidth]{0.4\textwidth}{12.7-17.png}

(1) 如图，①若 $\angle B = 30^\circ$，$\angle C = 40^\circ$，求 $\angle EAF$ 的度数；

②如果 $BC = 10$，求 $\triangle EAF$ 的周长；

③若 $AE \perp AF$，则 $\angle BAC = \underline{\quad }^\circ$。

(2) 若 $\angle BAC = n^\circ$，则 $\angle EAF = \underline{\quad }^\circ$。（用含 $n$ 代数式表示）

\vspace{7cm}


\item 如图，在 $\triangle ABC$ 中，$AB$ 边的垂直平分线 $l_1$ 交 $BC$ 于点 $D$，$AC$ 边的垂直平分线 $l_2$ 交 $BC$ 于点 $E$，$l_1$ 与 $l_2$ 相交于点 $O$，连接 $OB$、$OC$，若 $\triangle ADE$ 的周长为 $6\text{cm}$，$\triangle OBC$ 的周长为 $16\text{cm}$。

% 插入图片：此处应有第18题图

1. 求线段 $BC$ 的长；

2. 连接 $OA$，求线段 $OA$ 的长；

3. 若 $\angle BAC = 120^\circ$，求 $\angle DAE$ 的度数。

\rightquestionimage[width=0.4\textwidth]{0.4\textwidth}{12.7-18.png}


\end{enumerate}
